\chapter{Некоторые общематематические понятия и обозначения}

% -----------------------------------------------------------------------------
\lecture{Логическая символика}

\subsection{Связки и скобки}

\begin{definition}{Логические символы}{logical-symbols}
  Наряду с этими специальными символами, которые будут вводиться по мере
  надобности, мы используем распространенные символы математической логики
  $\neg$, $\land$, $\lor$, $\Rightarrow$, $\Leftrightarrow$ для обозначения
  соответственно отрицания «не» и связок «и», «или», «влечет»,
  «равносильно».
\end{definition}

\begin{properties}{Порядок приоритета символов}{logical-priority}
  Условимся с этой целью о следующем порядке приоритета символов:
  \[
    \neg,\quad \land,\quad \lor,\quad \Rightarrow,\quad \Leftrightarrow.
  \]
\end{properties}

\begin{definition}{Необходимое и достаточное условия}{necessary-sufficient}
  Записи $A\Rightarrow B$, означающей, что $A$ влечет $B$ или, что то же
  самое, $B$ следует из $A$, мы часто будем придавать другую словесную
  интерпретацию, говоря, что $B$ есть необходимый признак или необходимое
  условие $A$ и, в свою очередь, $A$ — достаточное условие или достаточный
  признак $B$.

  Итак, запись $A\Leftrightarrow B$ означает, что $A$ влечет $B$ и,
  одновременно, $B$ влечет $A$.
\end{definition}

\begin{definition}{Неразделительное «или»}{inclusive-or}
  Следует, однако, обратить внимание на то, что в выражении $A\lor B$ союз
  или неразделительный, т.~е. высказывание $A\lor B$ считается верным, если
  истинно хотя бы одно из высказываний $A$, $B$.
\end{definition}

\begin{properties}{Принцип исключенного третьего}{excluded-middle}
  При доказательстве от противного мы будем использовать также принцип
  исключенного третьего, в силу которого высказывание $A\lor\neg A$ ($A$
  или не $A$) считается истинным независимо от конкретного содержания
  высказывания $A$. Следовательно, мы одновременно принимаем, что
  $\neg(\neg A)\Leftrightarrow A$, т.~е. повторное отрицание равносильно
  исходному высказыванию.
\end{properties}

\begin{definition}{Равенство по определению}{definition-equality}
  Условимся также, когда это будет удобно, вводить определения посредством
  специального символа $:=$ (равенство по определению), в котором двоеточие
  ставится со стороны определяемого объекта.
\end{definition}

% -----------------------------------------------------------------------------
\lecture{Множества и элементарные операции над множествами}

\subsection{Понятие множества}

\begin{definition}{Элементы множества}{set-elements}
  Если $x$ — объект, $P$ — свойство, $P(x)$ — обозначение того, что $x$
  обладает свойством $P$, то через $\{x\mid P(x)\}$ обозначают весь класс
  объектов, обладающих свойством $P$. Объекты, составляющие класс или
  множество, называют элементами класса или множества.
\end{definition}

\subsection{Отношение включения}

\begin{definition}{Равенство множеств}{set-equality}
  Два множества равны, когда они состоят из одних и тех же элементов.
\end{definition}

\begin{definition}{Включение множеств}{set-inclusion}
  Если любой элемент множества $A$ является элементом множества $B$, то
  пишут $A\subset B$ или $B\supset A$ и говорят, что множество $A$ является
  подмножеством множества $B$, или что $B$ содержит $A$, или что $B$ включает
  в себя $A$. В связи с этим отношение $A\subset B$ между множествами $A$,
  $B$ называется отношением включения.
  \[
    (A\subset B):=\forall x\bigl((x\in A)\Rightarrow(x\in B)\bigr).
  \]
\end{definition}

\begin{definition}{Строгое включение}{strict-inclusion}
  Если $A\subset B$ и $A\ne B$, то будем говорить, что включение
  $A\subset B$ строгое или что $A$ — собственное подмножество $B$.
\end{definition}

\begin{properties}{Критерий равенства множеств}{set-equality-criterion}
  \[
    (A=B)\Leftrightarrow(A\subset B)\land(B\subset A).
  \]
\end{properties}

\begin{definition}{Пустое множество}{empty-set}
  Если в качестве $P$ взять свойство, которым не обладает ни один элемент
  множества $M$, например $P(x):=(x\ne x)$, то мы получим множество
  \[
    \varnothing=\{x\in M\mid x\ne x\},
  \]
  называемое пустым подмножеством множества $M$.
\end{definition}

\subsection{Простейшие операции над множествами}

\begin{definition}{Объединение множеств}{set-union}
  Объединением множеств $A$ и $B$ называется множество
  \[
    A\cup B:=\{x\in M\mid(x\in A)\lor(x\in B)\},
  \]
  состоящее из тех и только тех элементов множества $M$, которые содержатся
  хотя бы в одном из множеств $A$, $B$.
\end{definition}

\begin{definition}{Пересечение множеств}{set-intersection}
  Пересечением множеств $A$ и $B$ называется множество
  \[
    A\cap B:=\{x\in M\mid(x\in A)\land(x\in B)\},
  \]
  образованное теми и только теми элементами множества $M$, которые
  принадлежат одновременно множествам $A$ и $B$.
\end{definition}

\begin{definition}{Разность множеств}{set-difference}
  Разностью между множеством $A$ и множеством $B$ называется множество
  \[
    A\setminus B:=\{x\in M\mid(x\in A)\land(x\notin B)\},
  \]
  состоящее из тех элементов множества $A$, которые не содержатся в
  множестве $B$.
\end{definition}

\begin{definition}{Дополнение множества}{set-complement}
  Разность между множеством $M$ и содержащимся в нем подмножеством $A$
  обычно называют дополнением $A$ в $M$ и обозначают через $C_MA$ или $CA$,
  когда из контекста ясно, в каком множестве ищется дополнение к $A$.
\end{definition}

\begin{properties}{Правила де Моргана}{de-morgan-sets}
  \begin{align}
    C_M(A\cup B)&=C_MA\cap C_MB,\\
    C_M(A\cap B)&=C_MA\cup C_MB.
  \end{align}
\end{properties}

\subsection{Прямое произведение множеств}

\begin{definition}{Неупорядоченная и упорядоченная пары}{ordered-pair}
  Для любых двух множеств $A$, $B$ можно образовать новое множество — пару
  $\{A,B\}=\{B,A\}$, элементами которого являются множества $A$ и $B$ и
  только они. Это множество состоит из двух элементов, если $A\ne B$, и из
  одного элемента, если $A=B$.

  Указанное множество называют неупорядоченной парой множеств $A$, $B$, в
  отличие от упорядоченной пары $(A,B)$, в которой элементы $A$, $B$ наделены
  дополнительными признаками, выделяющими первый и второй элементы пары
  $\{A,B\}$. Равенство
  \[
    (A,B)=(C,D)
  \]
  упорядоченных пар по определению означает, что $A=C$ и $B=D$.
\end{definition}

\begin{definition}{Прямое произведение множеств}{cartesian-product}
  Пусть теперь $X$ и $Y$ — произвольные множества. Множество
  \[
    X\times Y:=\{(x,y)\mid(x\in X)\land(y\in Y)\},
  \]
  образованное всеми упорядоченными парами $(x,y)$, первый член которых есть
  элемент из $X$, а второй член — элемент из $Y$, называется прямым или
  декартовым произведением множеств $X$ и $Y$ (в таком порядке!).
\end{definition}

\begin{definition}{Проекции упорядоченной пары}{pair-projections}
  В упорядоченной паре $z=(x_1,x_2)$, являющейся элементом прямого
  произведения $Z=X_1\times X_2$ множеств $X_1$ и $X_2$, элемент $x_1$
  называется первой проекцией пары $z$ и обозначается через $\operatorname{pr}_1z$,
  а элемент $x_2$ — второй проекцией пары $z$ и обозначается через
  $\operatorname{pr}_2z$.

  Проекции упорядоченной пары по аналогии с терминологией аналитической
  геометрии часто называют (первой и второй) координатами пары.
\end{definition}

% -----------------------------------------------------------------------------
\lecture{Функция}

\subsection{Понятие функции (отображения)}

\begin{definition}{Функция}{function}
  Пусть $X$ и $Y$ — какие-то множества.

  Говорят, что имеется функция, определенная на $X$ со значениями в $Y$, если
  в силу некоторого закона $f$ каждому элементу $x\in X$ соответствует элемент
  $y\in Y$.

  В этом случае множество $X$ называется областью определения функции;
  символ $x$ его общего элемента — аргументом функции или независимой
  переменной; соответствующий конкретному значению $x_0\in X$ аргумента $x$
  элемент $y_0\in Y$ называют значением функции на элементе $x_0$ или
  значением функции при значении аргумента $x=x_0$ и обозначают через
  $f(x_0)$.
\end{definition}

\begin{definition}{Множество значений функции}{function-range}
  Множество
  \[
    f(X):=\{y\in Y\mid\exists x\,((x\in X)\land(y=f(x)))\}
  \]
  всех значений функции, которые она принимает на элементах множества $X$,
  будем называть множеством значений или областью значений функции.
\end{definition}

\begin{definition}{Равенство функций}{function-equality}
  Две функции $f_1$, $f_2$ считаются совпадающими или равными, если они имеют
  одну и ту же область определения $X$ и на любом элементе $x\in X$ значения
  $f_1(x)$, $f_2(x)$ этих функций совпадают. В этом случае пишут $f_1=f_2$.
\end{definition}

\begin{definition}{Сужение и продолжение функции}{restriction-extension}
  Если $A\subset X$, а $f:X\to Y$ — некоторая функция, то через $f|_A$ или
  $f|A$ обозначают функцию $\varphi:A\to Y$, совпадающую с $f$ на множестве
  $A$, т.~е. $f|_A(x)=\varphi(x)=f(x)$, если $x\in A$. Функция $f|_A$
  называется сужением или ограничением функции $f$ на множество $A$, а
  функция $f:X\to Y$ по отношению к функции $\varphi=f|_A:A\to Y$
  называется распространением или продолжением функции $\varphi$ на
  множество $X$.
\end{definition}

\begin{definition}{Области отправления и прибытия}{departure-arrival}
  Для обозначения любого множества $X$, содержащего область определения
  функции, иногда используется термин область отправления функции, а любое
  множество $Y$, содержащее область значений функции, называют тогда
  областью ее прибытия.
\end{definition}

\subsection{Простейшая классификация отображений}

\begin{definition}{Образ множества}{set-image}
  Образом множества $A\subset X$ при отображении $f:X\to Y$ называют
  множество
  \[
    f(A):=\{y\in Y\mid\exists x\,((x\in A)\land(y=f(x)))\}
  \]
  тех элементов $Y$, которые являются образами элементов множества $A$.
\end{definition}

\begin{definition}{Прообраз множества}{set-preimage}
  Множество
  \[
    f^{-1}(B):=\{x\in X\mid f(x)\in B\}
  \]
  тех элементов $X$, образы которых содержатся в $B$, называют прообразом
  (или полным прообразом) множества $B\subset Y$.
\end{definition}

\begin{definition}{Сюръекция, инъекция и биекция}{mapping-classes}
  Про отображение $f:X\to Y$ говорят, что оно
  \begin{itemize}
    \item сюръективно (или есть отображение $X$ на $Y$), если $f(X)=Y$;
    \item инъективно (или есть вложение, инъекция), если для любых элементов
      $x_1$, $x_2$ множества $X$
      \[
        (f(x_1)=f(x_2))\Rightarrow(x_1=x_2),
      \]
      т.~е. различные элементы имеют различные образы;
    \item биективно (или взаимно однозначно), если оно сюръективно и
      инъективно одновременно.
  \end{itemize}
\end{definition}

\begin{definition}{Обратное отображение}{inverse-map}
  Если отображение $f:X\to Y$ биективно, т.~е. является взаимно однозначным
  соответствием между элементами множеств $X$ и $Y$, то естественно возникает
  отображение
  \[
    f^{-1}:Y\to X,
  \]
  которое определяется следующим образом: если $f(x)=y$, то $f^{-1}(y)=x$,
  т.~е. элементу $y\in Y$ ставится в соответствие тот элемент $x\in X$,
  образом которого при отображении $f$ является $y$.

  Это отображение называют обратным по отношению к исходному отображению
  $f$.
\end{definition}

\begin{properties}{Свойства обратного отображения}{inverse-properties}
  Из построения обратного отображения видно, что $f^{-1}:Y\to X$ само
  является биективным и что обратное к нему отображение
  $(f^{-1})^{-1}:X\to Y$ совпадает с $f:X\to Y$.

  Таким образом, свойство двух отображений быть обратными является
  взаимным: если $f^{-1}$ — обратное для $f$, то, в свою очередь, $f$ —
  обратное для $f^{-1}$.
\end{properties}

\subsection{Композиция функций и взаимно обратные отображения}

\begin{definition}{Композиция отображений}{map-composition}
  Если отображения $f:X\to Y$ и $g:Y\to Z$ таковы, что одно из них (в нашем
  случае $g$) определено на множестве значений другого ($f$), то можно
  построить новое отображение
  \[
    g\circ f:X\to Z,
  \]
  значения которого на элементах множества $X$ определяются формулой
  \[
    (g\circ f)(x):=g(f(x)).
  \]
  Построенное составное отображение $g\circ f$ называют композицией
  отображения $f$ и отображения $g$ (в таком порядке!).
\end{definition}

\begin{properties}{Ассоциативность композиции}{composition-associativity}
  Операцию композиции иногда приходится проводить несколько раз подряд, и
  в этой связи полезно отметить, что она ассоциативна, т.~е.
  \[
    h\circ(g\circ f)=(h\circ g)\circ f.
  \]
\end{properties}

\begin{definition}{Тождественное отображение}{identity-map}
  Отображение $f:X\to X$, сопоставляющее каждому элементу множества $X$ его
  самого, т.~е. $x\mapsto x$, будем обозначать через $e_X$ и называть
  тождественным отображением множества $X$.
\end{definition}

\begin{lemma}{О композиции отображений}{composition-lemma}
  \[
    (g\circ f=e_X)\Rightarrow
    (g\text{ сюръективно})\land(f\text{ инъективно}).
  \]
\end{lemma}

\begin{proposition}{О взаимно обратных отображениях}{inverse-maps}
  Отображения $f:X\to Y$, $g:Y\to X$ являются биективными и взаимно
  обратными в том и только в том случае, когда $g\circ f=e_X$ и
  $f\circ g=e_Y$.
\end{proposition}

\subsection{Функция как отношение. График функции}

\begin{definition}{Отношение}{relation}
  Отношением $R$ называют любое множество упорядоченных пар $(x,y)$.

  Множество $X$ первых элементов упорядоченных пар, составляющих $R$,
  называют областью определения отношения $R$, а множество $Y$ вторых
  элементов этих пар — областью значений отношения $R$.
\end{definition}

\begin{definition}{Области отправления и прибытия отношения}{relation-domains}
  Любое множество, содержащее область определения отношения, называют
  областью отправления этого отношения. Множество, содержащее область
  значений отношения, называют областью прибытия отношения.
\end{definition}

\begin{definition}{Отношение эквивалентности}{equivalence-relation}
  Любое отношение $R$, обладающее перечисленными тремя свойствами, т.~е.
  рефлексивное, симметричное и транзитивное, принято называть отношением
  эквивалентности.
  \begin{align*}
    aRa &\quad\text{(рефлексивность)};\\
    aRb&\Rightarrow bRa \quad\text{(симметричность)};\\
    (aRb)\land(bRc)&\Rightarrow aRc \quad\text{(транзитивность)}.
  \end{align*}
\end{definition}

\begin{definition}{Отношение частичного порядка}{partial-order}
  Отношение между парами элементов некоторого множества $X$, обладающее
  указанными тремя свойствами, принято называть отношением частичного
  порядка на множестве $X$.
  \begin{align*}
    aRa &\quad\text{(рефлексивность)};\\
    (aRb)\land(bRc)&\Rightarrow aRc \quad\text{(транзитивность)};\\
    (aRb)\land(bRa)&\Rightarrow a\mathbin{\Delta}b,
      \text{ т.~е. }a=b \quad\text{(антисимметричность)}.
  \end{align*}
\end{definition}

\begin{definition}{Линейно упорядоченное множество}{linear-order}
  Если кроме отмеченных двух свойств, определяющих отношение частичного
  порядка, выполнено условие, что
  \[
    \forall a\,\forall b\,((aRb)\lor(bRa)),
  \]
  т.~е. любые два элемента множества $X$ сравнимы, то отношение $R$
  называется отношением порядка, а множество $X$ с определенным на нем
  отношением порядка называется линейно упорядоченным.
\end{definition}

\begin{definition}{Функциональное отношение}{functional-relation}
  Отношение $R$ называется функциональным, если
  \[
    (xRy_1)\land(xRy_2)\Rightarrow(y_1=y_2).
  \]
  Функциональное отношение называют функцией.
\end{definition}

\begin{definition}{График функции}{function-graph}
  Графиком функции $f:X\to Y$, понимаемой в смысле исходного описания,
  называют подмножество $\Gamma$ прямого произведения $X\times Y$, элементы
  которого имеют вид $(x,f(x))$. Итак,
  \[
    \Gamma:=\{(x,y)\in X\times Y\mid y=f(x)\}.
  \]
\end{definition}

% -----------------------------------------------------------------------------
\lecture{Некоторые дополнения}

\subsection{Мощность множества (кардинальные числа)}

\begin{definition}{Равномощные множества}{equinumerous-sets}
  Говорят, что множество $X$ равномощно множеству $Y$, если существует
  биективное отображение $X$ на $Y$, т.~е. каждому элементу $x\in X$
  сопоставляется элемент $y\in Y$, причем различным элементам множества $X$
  отвечают различные элементы множества $Y$ и каждый элемент $y\in Y$
  сопоставлен некоторому элементу множества $X$.
\end{definition}

\begin{definition}{Мощность множества}{cardinality}
  Класс, которому принадлежит множество $X$, называется мощностью множества
  $X$, а также кардиналом или кардинальным числом множества $X$ и
  обозначается символом $\card X$. Если $X\sim Y$, то пишут
  $\card X=\card Y$.
\end{definition}

\begin{definition}{Сравнение кардинальных чисел}{cardinal-comparison}
  Говорят, что кардинальное число множества $X$ не больше кардинального
  числа множества $Y$, и пишут $\card X\le\card Y$, если $X$ равномощно
  некоторому подмножеству множества $Y$.

  Итак,
  \[
    (\card X\le\card Y):=(\exists Z\subset Y\mid\card X=\card Z).
  \]
\end{definition}

\begin{definition}{Конечные и бесконечные множества}{finite-infinite}
  Множество называется конечным (по Дедекинду), если оно не равномощно
  никакому своему собственному подмножеству; в противном случае оно
  называется бесконечным.
\end{definition}

\begin{properties}{Свойства сравнения мощностей}{cardinal-properties}
  \begin{enumerate}[label=\arabic*${}^{\circ}$]
    \item $(\card X\le\card Y)\land(\card Y\le\card Z)
      \Rightarrow(\card X\le\card Z)$ (очевидно);\par
    \item $(\card X\le\card Y)\land(\card Y\le\card X)
      \Rightarrow(\card X=\card Y)$ (теорема Шрёдера—Бернштейна);\par
    \item $\forall X\,\forall Y\,(\card X\le\card Y)
      \lor(\card Y\le\card X)$ (теорема Кантора).
  \end{enumerate}
\end{properties}

\begin{definition}{Строгое сравнение мощностей}{strict-cardinal-comparison}
  Говорят, что мощность множества $X$ меньше мощности множества $Y$, и пишут
  $\card X<\card Y$, если $\card X\le\card Y$ и в то же время
  $\card X\ne\card Y$. Итак,
  \[
    (\card X<\card Y):=(\card X\le\card Y)\land(\card X\ne\card Y).
  \]
\end{definition}

\begin{theorem}{Кантора}{cantor}
  \[
    \card X<\card\powerset{X}.
  \]
\end{theorem}

\begin{corollary}{О мощностях бесконечных множеств}{different-infinities}
  Эта теорема, в частности, показывает, что если бесконечные множества
  существуют, то и «бесконечности» бывают разные.
\end{corollary}

\subsection{Об аксиоматике теории множеств}

\begin{axiom}{Объемности}{set-extensionality}
  Множества $A$ и $B$ равны тогда и только тогда, когда они имеют одни и те
  же элементы.
\end{axiom}

\begin{axiom}{Выделения}{set-separation}
  Любому множеству $A$ и свойству $P$ отвечает множество $B$, элементы
  которого суть те и только те элементы множества $A$, которые обладают
  свойством $P$.
\end{axiom}

\begin{corollary}{Единственность пустого множества}{unique-empty-set}
  Из аксиомы выделения следует, что существует пустое подмножество
  $\varnothing_X=\{x\in X\mid x\ne x\}$ в любом множестве $X$, а с учетом
  аксиомы объемности заключаем, что для любых множеств $X$ и $Y$ выполнено
  $\varnothing_X=\varnothing_Y$, т.~е. пустое множество единственно. Его
  обозначают символом $\varnothing$.
\end{corollary}

\begin{axiom}{Объединения}{set-union-axiom}
  Для любого множества $M$ множеств существует множество $\bigcup M$,
  называемое объединением множества $M$, состоящее из тех и только тех
  элементов, которые содержатся в элементах множества $M$.
\end{axiom}

\begin{axiom}{Пары}{set-pairing}
  Для любых множеств $X$ и $Y$ существует множество $Z$ такое, что $X$ и
  $Y$ являются его единственными элементами.
\end{axiom}

\begin{axiom}{Множества подмножеств}{power-set-axiom}
  Для любого множества $X$ существует множество $\powerset{X}$, состоящее
  из тех и только тех элементов, которые являются подмножествами множества
  $X$.
\end{axiom}

\begin{definition}{Последователь множества}{set-successor}
  Для того чтобы сформулировать следующую аксиому, введем понятие
  последователя $X^+$ множества $X$. Положим по определению
  $X^+=X\cup\{X\}$. Короче, к $X$ добавлено одноэлементное множество
  $\{X\}$.
\end{definition}

\begin{definition}{Индуктивное множество}{inductive-set}
  Далее, множество назовем индуктивным, если оно содержит в качестве
  элемента пустое множество и последователь любого своего элемента.
\end{definition}

\begin{axiom}{Бесконечности}{infinity-axiom}
  Индуктивные множества существуют.
\end{axiom}

\begin{axiom}{Подстановки}{replacement-axiom}
  Пусть $F(x,y)$ — такое высказывание (точнее, формула), что при любом
  $x_0$ из множества $X$ существует и притом единственный объект $y_0$
  такой, что $F(x_0,y_0)$ истинно. Тогда объекты $y$, для каждого из которых
  существует элемент $x\in X$ такой, что $F(x,y)$ истинно, образуют
  множество.
\end{axiom}

\begin{axiom}{Выбора}{choice-axiom}
  Для любого семейства непустых попарно непересекающихся множеств
  существует множество $C$ такое, что, каково бы ни было множество $X$
  данного семейства, множество $X\cap C$ состоит из одного элемента.
\end{axiom}

\subsection{Структура математических высказываний}

\begin{definition}{Квантор существования и единственности}{unique-existence}
  Запись
  \[
    \exists x\bigl(P(x)\land(\forall y\,(P(y)\Rightarrow y=x))\bigr)
  \]
  означает, что найдется объект $x$, обладающий свойством $P$ и такой, что
  если $y$ — любой объект, обладающий свойством $P$, то $y=x$. Короче:
  существует и притом единственный объект $x$, обладающий свойством $P$.
  Обычно это высказывание обозначают в виде $\exists!x\,P(x)$, и мы будем
  использовать такое сокращение.
\end{definition}

\begin{properties}{Сокращения для кванторов}{quantifier-shortcuts}
  \begin{align*}
    (\forall x\in X)\,P&:=\forall x\,(x\in X\Rightarrow P(x)),\\
    (\exists x\in X)\,P&:=\exists x\,(x\in X\land P(x)),\\
    (\forall x>a)\,P&:=\forall x\,(x\in\R\land x>a\Rightarrow P(x)),\\
    (\exists x>a)\,P&:=\exists x\,(x\in\R\land x>a\land P(x)).
  \end{align*}
\end{properties}

\begin{properties}{Отрицания высказываний}{logical-negations}
  Отрицание к высказыванию «для некоторого $x$ истинно $P(x)$» означает,
  что «для любого $x$ неверно $P(x)$», а отрицание к высказыванию «для любого
  $x$ истинно $P(x)$» означает, что «найдется $x$, что неверно $P(x)$».

  Итак,
  \begin{align*}
    \neg\exists x\,P(x)&\Leftrightarrow\forall x\,\neg P(x),\\
    \neg\forall x\,P(x)&\Leftrightarrow\exists x\,\neg P(x),\\
    \neg(P\land Q)&\Leftrightarrow\neg P\lor\neg Q,\\
    \neg(P\lor Q)&\Leftrightarrow\neg P\land\neg Q,\\
    \neg(P\Rightarrow Q)&\Leftrightarrow P\land\neg Q.
  \end{align*}
\end{properties}
